\section{Sprzężenia funktorów}\label{sec:adjunctions}

Sprzężenie opisuje sytuację, w której konstrukcja nowego obiektu jest \emph{najswobodniejsza} lub \emph{najbardziej uniwersalna} względem zapomnienia pewnej struktury. To jeden z najczęściej występujących wzorców we współczesnej algebrze i topologii.

\subsection{Definicja przez naturalne bijekcje}

\begin{definicja}[Para funktorów sprzężonych]\label{def:adjunction}
Niech $\mathcal C,\mathcal D$ będą lokalnie małymi kategoriami, a $F:\mathcal C\to\mathcal D$ i $G:\mathcal D\to\mathcal C$ funktorami. Mówimy, że $F$ jest \emph{lewym sprzężeniem} $G$, a $G$ \emph{prawym sprzężeniem} $F$, jeśli dla wszystkich $A\in\Ob(\mathcal C)$ i $B\in\Ob(\mathcal D)$ mamy bijekcje
\[
\Phi_{A,B}:\Hom_{\mathcal D}(F(A),B)\ \overset{\cong}{\longrightarrow}\ \Hom_{\mathcal C}(A,G(B)),
\]
naturalne jednocześnie względem $A$ i $B$. Piszemy $F\dashv G$.
\end{definicja}

Naturalność względem obu zmiennych oznacza konkretnie, że dla $u:A'\to A$, $v:B\to B'$ i $h:F(A)\to B$ zachodzą
$\Phi_{A',B}(h\circ F(u))=\Phi_{A,B}(h)\circ u$ oraz $\Phi_{A,B'}(v\circ h)=G(v)\circ\Phi_{A,B}(h)$.

\begin{przyklad}[Wolna przestrzeń wektorowa]\label{przyklad:free-vect}
Dla zbioru $S$ niech $k^{(S)}$ oznacza przestrzeń wektorową funkcji $a:S\to k$ o \emph{skończonym nośniku} (tj. zbiór $\{s:a(s)\neq0\}$ jest skończony). Dla $s\in S$ oznaczmy przez $e_s$ funkcję przyjmującą wartość $1$ w $s$ i $0$ poza $s$. Każdy wektor ma jednoznaczny zapis $\sum_{s\in S}a_s e_s$ z tylko skończenie wieloma niezerowymi współczynnikami. Definiujemy funktor $F:\Set\to\Vect_k$ przez $F(S)=k^{(S)}$ i $F(f)(e_s)=e_{f(s)}$ dla funkcji $f:S\to T$, a następnie rozszerzamy liniowo.

Jeśli $U:\Vect_k\to\Set$ jest funktorem zapominania, to każda funkcja $\varphi:S\to U(V)$ ma dokładnie jedno przedłużenie liniowe $\widetilde\varphi:k^{(S)}\to V$, dane przez $\widetilde\varphi(\sum a_s e_s)=\sum a_s\varphi(s)$. Odwrotnie, przekształceniu liniowemu $L:k^{(S)}\to V$ odpowiada funkcja $s\mapsto L(e_s)$. Otrzymujemy naturalną bijekcję
\[
\Hom_{\Vect_k}(k^{(S)},V)\cong\Hom_{\Set}(S,U(V)),
\]
czyli $F\dashv U$. Wolna przestrzeń jest tak skonstruowana, aby \emph{każda} funkcja ze zbioru generatorów jednoznacznie przedłużała się liniowo.
\end{przyklad}

\begin{przyklad}[Uniwersalny element wolnej przestrzeni]
Dla ustalonego zbioru $S$ i funktora $G:\Vect_k\to\Set$ określonego przez $G(V)=\Hom_{\Set}(S,U(V))$ reprezentującym obiektem jest $k^{(S)}$. Element uniwersalny $u\in G(k^{(S)})$ to funkcja $s\mapsto e_s$. Dla dowolnej $\varphi:S\to U(V)$ istnieje dokładnie jedno odwzorowanie liniowe $L:k^{(S)}\to V$, dla którego $G(L)(u)=\varphi$, czyli $L(e_s)=\varphi(s)$ dla wszystkich $s$. Jest to element uniwersalny odpowiadający sprzężeniu wolnej przestrzeni z funktorem zapominania.
\end{przyklad}

\begin{przyklad}[Grupa wolna]
Dla zbioru $S$ istnieje grupa wolna $F_{\Grp}(S)$ na generatorach $S$ (elementami są zredukowane słowa w symbolach $s$ oraz $s^{-1}$). Każdą funkcję $S\to U(G)$ do zbioru elementów grupy $G$ przedłużamy jednoznacznie do homomorfizmu $F_{\Grp}(S)\to G$. Otrzymujemy sprzężenie $F_{\Grp}:\Set\rightleftarrows\Grp:U$. Dla $S=\{*\}$ grupa wolna jest izomorficzna z $\mathbb Z$, co zgadza się z przykładem w sekcji~\ref{sec:yoneda}.
\end{przyklad}

\begin{przyklad}[Dwie konstrukcje topologiczne]
Niech $U:\Top\to\Set$ zapomina topologię. Dla zbioru $S$ oznaczmy przez $D(S)$ zbiór $S$ z topologią dyskretną (wszystkie podzbiory są otwarte), a przez $I(S)$ zbiór $S$ z topologią antydyskretną (otwarte są tylko $\varnothing$ i $S$). Każda funkcja $D(S)\to X$ jest ciągła, a każda funkcja $X\to I(S)$ jest ciągła. Dlatego mamy dwa sprzężenia
\[
D\dashv U\dashv I.
\]
Rzeczywiście $\Hom_{\Top}(D(S),X)\cong\Hom_{\Set}(S,U(X))$ oraz $\Hom_{\Set}(U(X),S)\cong\Hom_{\Top}(X,I(S))$. Jeden funktor zapominania może zatem mieć zarówno lewe, jak i prawe sprzężenie.
\end{przyklad}

\begin{przyklad}[Przekątna, iloczyn i koprodukt]\label{przyklad:diag-adj}
W $\Set$ niech $\Delta:\Set\to\Set\times\Set$ będzie funktorem przekątnym, $\Delta(X)=(X,X)$. Niech $\times:\Set\times\Set\to\Set$ oznacza iloczyn dwóch zbiorów, a $\sqcup:\Set\times\Set\to\Set$ ich sumę rozłączną. Uniwersalne własności iloczynu i koproduktu można zapisać jako dwa sprzężenia
\[
\sqcup\ \dashv\ \Delta\ \dashv\ \times.
\]
Pierwsze wynika z $\Hom(A\sqcup B,X)\cong\Hom(A,X)\times\Hom(B,X)$, a drugie z $\Hom(X,A\times B)\cong\Hom(X,A)\times\Hom(X,B)$.
\end{przyklad}


\begin{przyklad}[Sprzężenie funkcji monotonicznych na liczbach całkowitych]\label{przyklad:adj-poset}
Traktujmy $(\mathbb Z,\leq)$ jako kategorię. Funkcje $F,G:\mathbb Z\to\mathbb Z$ dane przez $F(n)=2n$ i $G(m)=\lfloor m/2\rfloor$ są monotoniczne, a więc są funktorami. Dla dowolnych liczb całkowitych $n,m$ zachodzi
$F(n)\leq m\ \Longleftrightarrow\ 2n\leq m\ \Longleftrightarrow\ n\leq\lfloor m/2\rfloor\ \Longleftrightarrow\ n\leq G(m)$.
Każda ze stron oznacza niepustość odpowiedniego zbioru morfizmów, a skoro te zbiory są puste lub jednoelementowe, dostajemy bijekcję naturalną $\Hom(F(n),m)\cong\Hom(n,G(m))$. Zatem $F\dashv G$. Na przykład dla $n=3,m=7$ nierówności $6\leq7$ i $3\leq3$ są prawdziwe; dla $n=4,m=7$ obie są fałszywe.
\end{przyklad}

\begin{przyklad}[Obraz, przeciwobraz i kwantyfikator ``dla każdego'']\label{przyklad:adj-inverse-image}
Niech $f:X\to Y$ będzie funkcją. Na posetach podzbiorów mamy trzy monotoniczne funkcje:
\begin{itemize}
\item $\exists_f:\mathcal P(X)\to\mathcal P(Y)$, $\exists_f(A)=f(A)$;
\item $f^{-1}:\mathcal P(Y)\to\mathcal P(X)$, $f^{-1}(B)=\{x:f(x)\in B\}$;
\item $\forall_f:\mathcal P(X)\to\mathcal P(Y)$, $\forall_f(A)=\{y\in Y:f^{-1}(\{y\})\subseteq A\}$.
\end{itemize}
Dla $A\subseteq X$, $B\subseteq Y$ zachodzą równoważności
\[
f(A)\subseteq B\ \Longleftrightarrow\ A\subseteq f^{-1}(B),\qquad
f^{-1}(B)\subseteq A\ \Longleftrightarrow\ B\subseteq\forall_f(A).
\]
Są to dokładnie własności dwóch sprzężeń $\exists_f\dashv f^{-1}\dashv\forall_f$. Pierwszą równoważność sprawdzamy bezpośrednio: $f(x)\in B$ dla wszystkich $x\in A$ wtedy i tylko wtedy, gdy $x\in f^{-1}(B)$ dla wszystkich $x\in A$. Druga mówi, że każde włókno nad $y\in B$ jest zawarte w $A$ wtedy i tylko wtedy, gdy $y\in\forall_f(A)$.

Przykładowo dla $X=\{a,b,c\}$, $Y=\{0,1\}$, $f(a)=0$, $f(b)=f(c)=1$ i $A=\{a,b\}$ otrzymujemy $\exists_f(A)=\{0,1\}$, lecz $\forall_f(A)=\{0\}$. Warto uwzględnić także puste włókna: jeśli $y\notin f(X)$, to $y\notin\exists_f(A)$, lecz $y\in\forall_f(A)$ dla każdego $A\subseteq X$, gdyż $\varnothing\subseteq A$. Obie strony sprzężeń mają więc wyraźnie różne zadania: ``istnieje element włókna'' i ``wszystkie elementy włókna''.
\end{przyklad}

\begin{przyklad}[Naturalność sprzężenia wolnej przestrzeni]
Dla zbioru $S$, przestrzeni $V$, funkcji $\varphi:S\to U(V)$ odpowiadające jej liniowe przedłużenie spełnia $\widetilde\varphi(e_s)=\varphi(s)$. Sprawdźmy obie zmienne. Jeśli $a:S'\to S$, to $(\widetilde\varphi\circ F(a))(e_{s'})=\widetilde\varphi(e_{a(s')})=\varphi(a(s'))$, czyli liniowy morfizm odpowiadający $\varphi\circ a$. Jeśli $b:V\to W$ jest liniowe, to $(b\circ\widetilde\varphi)(e_s)=b(\varphi(s))$, a zatem odpowiada funkcji $U(b)\circ\varphi$. Wykazaliśmy naturalność bijekcji z przykładu~\ref{przyklad:free-vect} względem $S$ i $V$, a nie tylko bijektywność dla każdej ustalonej pary.
\end{przyklad}


\subsection{Jednostka, kojednostka i tożsamości trójkątne}

\begin{definicja}[Jednostka i kojednostka sprzężenia]\label{def:unit-counit}
Dla sprzężenia $F:\mathcal C\rightleftarrows\mathcal D:G$ z bijekcjami $\Phi_{A,B}$ definiujemy \emph{jednostkę} $\eta:\id_{\mathcal C}\Rightarrow GF$ przez $\eta_A:=\Phi_{A,F(A)}(\id_{F(A)})$. Definiujemy \emph{kojednostkę} $\varepsilon:FG\Rightarrow\id_{\mathcal D}$ przez $\varepsilon_B:=\Phi_{G(B),B}^{-1}(\id_{G(B)})$.
\end{definicja}

W sprzężeniach konstrukcji swobodnej z funktorem zapominania jednostka opisuje odwzorowanie generatorów do obiektu swobodnego, a kojednostka jego realizację w obiekcie z daną strukturą. W ogólnym sprzężeniu jednostka i kojednostka nie muszą być odpowiednio monomorfizmem i epimorfizmem. W przykładzie wolnej przestrzeni wektorowej $\eta_S(s)=e_s$, natomiast
$\varepsilon_V:k^{(U(V))}\to V$ spełnia $\varepsilon_V(e_v)=v$ i jest liniowa.

\begin{twierdzenie}[Tożsamości trójkątne i odtworzenie sprzężenia]\label{tw:triangles}
Dla każdego sprzężenia $F\dashv G$ jednostka $\eta$ i kojednostka $\varepsilon$ są transformacjami naturalnymi spełniającymi dla wszystkich $A\in\Ob(\mathcal C)$ i $B\in\Ob(\mathcal D)$ równości
\[
\varepsilon_{F(A)}\circ F(\eta_A)=\id_{F(A)},\qquad G(\varepsilon_B)\circ\eta_{G(B)}=\id_{G(B)}.
\]
Odwrotnie, jeśli dane są funktory $F,G$ i transformacje naturalne $\eta:\id_{\mathcal C}\Rightarrow GF$, $\varepsilon:FG\Rightarrow\id_{\mathcal D}$ spełniające te równości, to definiują sprzężenie $F\dashv G$.
\end{twierdzenie}
\begin{proof}
Najpierw wychodzimy od naturalnych bijekcji $\Phi_{A,B}$ z definicji sprzężenia. Zdefiniowane w sekcji jednostka $\eta$ i kojednostka $\varepsilon$ są naturalne: dla $u:A\to A'$ naturalność $\Phi$ w pierwszej i drugiej zmiennej daje
$GF(u)\eta_A=\Phi_{A,F(A')}(F(u))=\eta_{A'}u$. Dla $v:B\to B'$ naturalność bijekcji daje
\[
\begin{aligned}
\Phi_{G(B),B'}(v\circ\varepsilon_B)
&=G(v)\circ\Phi_{G(B),B}(\varepsilon_B)=G(v),\\
\Phi_{G(B),B'}(\varepsilon_{B'}\circ FG(v))
&=\Phi_{G(B'),B'}(\varepsilon_{B'})\circ G(v)=G(v).
\end{aligned}
\]
Injektywność $\Phi_{G(B),B'}$ daje $v\circ\varepsilon_B=\varepsilon_{B'}\circ FG(v)$, czyli naturalność kojednostki.

Z naturalności $\Phi$ w drugiej zmiennej, zastosowanej do $h:F(A)\to B$ i $\id_{F(A)}$, wynika ważny wzór
$\Phi_{A,B}(h)=G(h)\eta_A$.
Z naturalności $\Phi^{-1}$ w pierwszej zmiennej, zastosowanej do $u:A\to G(B)$ i $\id_{G(B)}$, wynika wzór odwrotny
$\Phi_{A,B}^{-1}(u)=\varepsilon_BF(u)$.
Zastosujmy drugi wzór przy $B=F(A)$ i $u=\eta_A$. Ponieważ $\eta_A=\Phi_{A,F(A)}(\id_{F(A)})$, mamy
\[
\varepsilon_{F(A)}F(\eta_A)
=\Phi_{A,F(A)}^{-1}(\eta_A)=\id_{F(A)}.
\]
Z kolei pierwszy wzór przy $A=G(B)$ i $h=\varepsilon_B$ daje
$G(\varepsilon_B)\eta_{G(B)}=\Phi_{G(B),B}(\varepsilon_B)=\id_{G(B)}$.
To dowodzi obu tożsamości trójkątnych, bez pomijania kroku prowadzącego od naturalności do wzorów na morfizmy sprzężone.

Odwrotnie, niech dane będą naturalne transformacje $\eta:\id_{\mathcal C}\Rightarrow GF$ i $\varepsilon:FG\Rightarrow\id_{\mathcal D}$ spełniające tożsamości trójkątne. Definiujemy dla $h:F(A)\to B$ oraz $u:A\to G(B)$ morfizmy
\[
\Phi_{A,B}(h)=G(h)\eta_A,\qquad
\Psi_{A,B}(u)=\varepsilon_BF(u).
\]
Mają one odpowiednie dziedziny i przeciwdziedziny. Naturalność $\varepsilon$ dla $h$ daje $\varepsilon_BFG(h)=h\varepsilon_{F(A)}$, stąd
$\Psi(\Phi(h))=\varepsilon_BFG(h)F(\eta_A)=h\varepsilon_{F(A)}F(\eta_A)=h$.
Naturalność $\eta$ dla $u$ daje $GF(u)\eta_A=\eta_{G(B)}u$, zatem
$\Phi(\Psi(u))=G(\varepsilon_B)GF(u)\eta_A=G(\varepsilon_B)\eta_{G(B)}u=u$.
Funkcje $\Phi_{A,B}$ i $\Psi_{A,B}$ są więc wzajemnie odwrotne. Naturalność $\Phi$ względem $u_0:A'\to A$ wynika z rachunku
$\Phi_{A',B}(hF(u_0))=G(h)GF(u_0)\eta_{A'}=G(h)\eta_Au_0=\Phi_{A,B}(h)u_0$.
Jej naturalność względem $v:B\to B'$ wynika z $\Phi_{A,B'}(vh)=G(v)G(h)\eta_A=G(v)\Phi_{A,B}(h)$. Tak otrzymujemy wymagane sprzężenie.
\end{proof}

\begin{przyklad}[Sprawdzenie pierwszego trójkąta dla wolnej przestrzeni]
Dla $S\in\Ob(\Set)$ i generatora $e_s\in F(S)$ mamy $F(\eta_S)(e_s)=e_{e_s}\in F(U(F(S)))$, po czym $\varepsilon_{F(S)}(e_{e_s})=e_s$. Obie strony $\varepsilon_{F(S)}\circ F(\eta_S)$ i $\id_{F(S)}$ są liniowe i zgadzają się na bazie $(e_s)_{s\in S}$, zatem są równe. To konkretna realizacja pierwszej tożsamości trójkątnej. Dla drugiej, przy $V\in\Ob(\Vect_k)$ i $v\in U(V)$ mamy $\eta_{U(V)}(v)=e_v$, a następnie $U(\varepsilon_V)(e_v)=v$. Złożenie $U(\varepsilon_V)\circ\eta_{U(V)}$ jest więc identycznością na $U(V)$.
\end{przyklad}


\begin{przyklad}[Oba trójkąty dla sprzężenia dyskretnego]
W sprzężeniu $D:\Set\rightleftarrows\Top:U$ z przykładu w sekcji~\ref{sec:adjunctions} mamy $UD=\id_{\Set}$ dosłownie, więc $\eta_S=\id_S$. Dla przestrzeni $X$ kojednostka $\varepsilon_X:D(U(X))\to X$ jest identycznością jako funkcja między zbiorami, a ciągłość wynika z dyskretności dziedziny. Pierwszy trójkąt mówi, że złożenie $D(S)\xrightarrow{D(\eta_S)}DUD(S)\xrightarrow{\varepsilon_{D(S)}}D(S)$ jest identycznością. Drugi mówi, że $U(X)\xrightarrow{\eta_{U(X)}}UDU(X)\xrightarrow{U(\varepsilon_X)}U(X)$ również jest identycznością. Oba wnioski są natychmiastowe po podstawieniu definicji, ale pokazują różnicę między identycznością funkcji a homeomorfizmem: $\varepsilon_X$ nie musi być homeomorfizmem.
\end{przyklad}


\subsection{Sprzężenia a zachowanie konstrukcji}

\begin{twierdzenie}[Lewe sprzężenia zachowują kogranice]\label{tw:adjoints-preserve}
Niech $F:\mathcal C\to\mathcal D$ i $G:\mathcal D\to\mathcal C$ będą funktorami ze sprzężeniem $F\dashv G$. Niech $J$ będzie małą kategorią i $D:J\to\mathcal C$ diagramem mającym kogranicę $(Q,(q_j)_j)$. Wówczas $(F(Q),(F(q_j))_j)$ jest kogranicą diagramu $F\circ D$ w $\mathcal D$. Dualnie, prawy sprzężony $G$ zachowuje wszystkie istniejące małe granice.
\end{twierdzenie}
\begin{proof}
Rodzina $F(q_j)$ jest kostożkiem, ponieważ dla $\alpha:j\to j'$ zachodzi $F(q_{j'})\circ F(D(\alpha))=F(q_{j'}\circ D(\alpha))=F(q_j)$.

Niech $(b_j:F(D(j))\to Y)_j$ będzie dowolnym kostożkiem. Definiujemy $a_j=\Phi_{D(j),Y}(b_j)$. Naturalność w pierwszej zmiennej daje
\[
a_{j'}\circ D(\alpha)
=\Phi_{D(j),Y}(b_{j'}\circ F(D(\alpha)))=a_j.
\]
Z własności uniwersalnej $Q$ istnieje dokładnie jeden $a:Q\to G(Y)$ taki, że $a\circ q_j=a_j$ dla każdego $j$. Niech $b=\Phi_{Q,Y}^{-1}(a)$. Ponownie z naturalności
\[
\Phi_{D(j),Y}(b\circ F(q_j))=a\circ q_j=a_j
=\Phi_{D(j),Y}(b_j).
\]
Injektywność daje $b\circ F(q_j)=b_j$. Jeśli $b'$ spełnia te same warunki, to $\Phi_{Q,Y}(b')\circ q_j=a_j$ dla wszystkich $j$, więc $\Phi_{Q,Y}(b')=a$ z jednoznaczności kogranicy, a następnie $b'=b$. Obejmuje to również pusty diagram.

Dla granic niech $E:J\to\mathcal D$ ma granicę $(L,(\ell_j)_j)$. Rodzina $G(\ell_j)$ jest stożkiem z funktorialności $G$. Dla dowolnego stożka $(c_j:X\to G(E(j)))_j$ definiujemy $d_j=\Phi_{X,E(j)}^{-1}(c_j)$. Dla $\alpha:j\to j'$ mamy
\[
\Phi_{X,E(j')}(E(\alpha)\circ d_j)
=G(E(\alpha))\circ c_j=c_{j'},
\]
więc $E(\alpha)\circ d_j=d_{j'}$. Istnieje zatem dokładnie jeden $d:F(X)\to L$ z $\ell_j\circ d=d_j$. Dla $c=\Phi_{X,L}(d)$ naturalność daje
\[
G(\ell_j)\circ c=\Phi_{X,E(j)}(\ell_j\circ d)=c_j.
\]
Jeśli $c'$ również spełnia te warunki, to $d'=\Phi_{X,L}^{-1}(c')$ ma $\ell_j\circ d'=d_j$, zatem $d'=d$ i $c'=c$. To dowodzi zachowania granicy przez $G$, również dla pustego diagramu.
\end{proof}

\begin{stwierdzenie}[Pełna wierność a jednostka i kojednostka]\label{cat:prop:adj-full-faithful}
W sprzężeniu $F\dashv G$ funktor $F$ jest w pełni wierny wtedy i tylko wtedy, gdy $\eta$ jest izomorfizmem naturalnym. Funktor $G$ jest w pełni wierny wtedy i tylko wtedy, gdy $\varepsilon$ jest izomorfizmem naturalnym.
\end{stwierdzenie}
\begin{proof}
Dla $A,A'\in\Ob(\mathcal C)$ złożenie odwzorowania na morfizmach $F$ z bijekcją sprzężenia wysyła $f:A\to A'$ na
\[
\Phi_{A,F(A')}(F(f))=GF(f)\circ\eta_A=\eta_{A'}\circ f.
\]
Jeżeli $\eta$ jest izomorfizmem, to ostatnie przyporządkowanie jest bijekcją, której odwrotność to składanie z $\eta_{A'}^{-1}$. Stąd $F$ jest w pełni wierny.

Odwrotnie, z pełnej wierności $F$ wynika, że dla każdego $A$ składanie z $\eta_{A'}$ jest bijekcją. Przy $A=GF(A')$ wybieramy jedyny $r:GF(A')\to A'$ taki, że $\eta_{A'}\circ r=\id_{GF(A')}$. Ponieważ
$\eta_{A'}\circ(r\circ\eta_{A'})=\eta_{A'}=\eta_{A'}\circ\id_{A'}$,
injektywność tej bijekcji przy $A=A'$ daje $r\circ\eta_{A'}=\id_{A'}$. Wszystkie składowe $\eta$ są izomorfizmami.

Dla $B,B'\in\Ob(\mathcal D)$ złożenie odwzorowania na morfizmach $G$ z odwrotną bijekcją sprzężenia wysyła $g:B\to B'$ na
\[
\Phi_{G(B),B'}^{-1}(G(g))=\varepsilon_{B'}\circ FG(g)=g\circ\varepsilon_B.
\]
Gdy $\varepsilon$ jest izomorfizmem, składanie z nią jest bijekcją, więc $G$ jest w pełni wierny. Jeżeli $G$ jest w pełni wierny, dla każdego $B'$ składanie z $\varepsilon_B$ jest bijekcją. Przy $B'=FG(B)$ otrzymujemy $t:B\to FG(B)$ z $t\circ\varepsilon_B=\id_{FG(B)}$. Ponieważ $(\varepsilon_B\circ t)\circ\varepsilon_B=\varepsilon_B=\id_B\circ\varepsilon_B$, injektywność przy $B'=B$ daje $\varepsilon_B\circ t=\id_B$. Zatem także $\varepsilon$ jest izomorfizmem naturalnym.
\end{proof}

\begin{przyklad}[Dlaczego wolna przestrzeń zachowuje koprodukty]
Dla zbiorów $S,T$ mamy kanoniczny izomorfizm $k^{(S\sqcup T)}\cong k^{(S)}\oplus k^{(T)}$. Można go sprawdzić bezpośrednio: funkcja o skończonym nośniku na sumie rozłącznej to para funkcji o skończonych nośnikach na jej składnikach. Twierdzenie~\ref{tw:adjoints-preserve} przewiduje ten izomorfizm bez rachunku na współrzędnych, ponieważ funktor wolnej przestrzeni jest lewym sprzężeniem i musi zachowywać koprodukty.
\end{przyklad}

\begin{uwaga}[Lewe i prawe są asymetryczne]
Lewe sprzężenie zachowuje kogranice, a prawe granice. Nie można bez dodatkowych założeń zamienić tych słów. Na przykład wolna przestrzeń $F:\Set\to\Vect_k$ nie zachowuje wszystkich iloczynów: dla zbiorów $S$ i $T$ o odpowiednio $2$ i $3$ elementach przestrzeń $F(S\times T)$ ma wymiar $6$, podczas gdy iloczyn $F(S)\times F(T)$ w $\Vect_k$ ma wymiar $2+3=5$.
\end{uwaga}


\begin{przyklad}[Prawy sprzężony zachowuje iloczyny]
Niech $U:\Grp\to\Set$ będzie funktorem zapominania. Ponieważ $U$ jest prawym sprzężeniem do funktora grupy wolnej, zachowuje granice. Dla grup $G,H$ dostajemy zatem kanoniczną bijekcję $U(G\times H)\cong U(G)\times U(H)$. Konkretna grupa $G\times H$ ma elementy $(g,h)$, więc bijekcja jest tu oczywista. Znaczenie twierdzenia polega jednak na tym, że działa ono także dla bardziej skomplikowanych granic, np. equalizatorów dowolnych homomorfizmów grup.
\end{przyklad}


